\documentclass{csexam}

\usepackage{listings}
\usepackage{xcolor}

% Python code styling
\lstset{
    language=Python,
    basicstyle=\ttfamily\small,
    keywordstyle=\color{blue}\bfseries,
    stringstyle=\color{red},
    commentstyle=\color{green!50!black}\itshape,
    showstringspaces=false,
    breaklines=true,
    frame=single,
    numbers=left,
    numberstyle=\tiny\color{gray}
}

% Course information
\course{CSCI 128: Coding for Problem Solving}
\courseshort{CSCI 128}
\exam{Quiz 6}
\examdate{Winter 2026}

\begin{document}

\maketitle

\begin{rules}
\textbf{Instructions:}
\begin{itemize}
    \item This quiz covers material from Lecture 6 (Combining Images and Collages)
    \item Answer all questions in the space provided
    \item Show your work where applicable
    \item Time limit: 15 minutes
\end{itemize}
\end{rules}

\vspace{0.2in}

\showgrade

\newpage

\begin{questions}

\question[2] What is the difference between a \textbf{canvas} and a \textbf{source image} in the context of image composition?

\vspace{1.5in}

\question[2] When pasting an image onto a canvas, the position argument specifies which corner of the pasted image?

\vspace{0.3cm}
\underline{\hspace{4in}}
\vspace{0.5in}

\question[3] You have a canvas of size 800$\times$600 and you paste a 200$\times$200 image at position $(700, 500)$.

\begin{parts}
    \part What is the right edge x-coordinate of the pasted image? \underline{\hspace{2in}}
    \part What is the bottom edge y-coordinate of the pasted image? \underline{\hspace{2in}}
    \part Does the image go off the canvas? Explain briefly.
    \vspace{0.8in}
\end{parts}

\question[3] For a grid collage with $N$ images arranged in a grid with \texttt{cols} columns, write the formulas to compute the row and column index of image $i$.

\vspace{0.3cm}
\texttt{row} $=$ \underline{\hspace{2.5in}}

\vspace{0.3cm}
\texttt{col} $=$ \underline{\hspace{2.5in}}

\vspace{0.3cm}
In a 3-column grid, what are the row and column of image 7?

\vspace{0.3cm}
\texttt{row} $=$ \underline{\hspace{1in}} \hspace{1cm} \texttt{col} $=$ \underline{\hspace{1in}}

\newpage

\question[3] What does the following code produce? Describe the output image in terms of its size and layout.

\begin{lstlisting}
from PIL import Image

img1 = Image.open("photo1.jpg").resize((300, 300))
img2 = Image.open("photo2.jpg").resize((300, 300))
img3 = Image.open("photo3.jpg").resize((300, 300))
img4 = Image.open("photo4.jpg").resize((300, 300))

canvas = Image.new("RGB", (600, 600))
canvas.paste(img1, (0, 0))
canvas.paste(img2, (300, 0))
canvas.paste(img3, (0, 300))
canvas.paste(img4, (300, 300))
\end{lstlisting}

\vspace{1.5in}

\question[3] In alpha blending, the result pixel is computed as:
\[
    \text{result} = \alpha \times \text{img1} + (1 - \alpha) \times \text{img2}
\]

\begin{parts}
    \part What does the output look like when $\alpha = 0.0$? \underline{\hspace{2.5in}}
    \part What does the output look like when $\alpha = 1.0$? \underline{\hspace{2.5in}}
    \part Two pixels are blended with $\alpha = 0.25$. Pixel from img1 is $(200, 100, 0)$ and from img2 is $(0, 100, 200)$. Calculate the resulting RGB values.

    \vspace{0.3cm}
    R $=$ \underline{\hspace{1in}} \hspace{0.5cm}
    G $=$ \underline{\hspace{1in}} \hspace{0.5cm}
    B $=$ \underline{\hspace{1in}}
\end{parts}

\newpage

\question[4] The function below implements a chromakey (green screen) effect. Read it carefully and answer the questions.

\begin{lstlisting}
def chromakey(foreground, background, key_color, threshold=50):
    fg_pixels = foreground.load()
    bg_pixels = background.load()
    result = Image.new("RGB", foreground.size)
    result_pixels = result.load()

    for x in range(foreground.size[0]):
        for y in range(foreground.size[1]):
            r, g, b = fg_pixels[x, y]
            distance = (abs(r - key_color[0]) + abs(g - key_color[1]) + abs(b - key_color[2]))
            if distance < threshold:
                result_pixels[x, y] = bg_pixels[x, y]
            else:
                result_pixels[x, y] = (r, g, b)
    return result
\end{lstlisting}

\begin{parts}
    \part What does the variable \texttt{distance} measure?
    \vspace{0.8in}

    \part A foreground pixel is $(10, 240, 20)$ and \texttt{key\_color} is $(0, 255, 0)$. Calculate \texttt{distance}.

    \vspace{0.3cm}
    \texttt{distance} $=$ \underline{\hspace{2in}}

    \part With \texttt{threshold = 50}, does this pixel get replaced by the background or kept from the foreground?

    \vspace{0.3cm}
    \underline{\hspace{4in}}
\end{parts}

\end{questions}

\end{document}
